\section{Discussion}
\label{Discussion}

Dataset-specific substitutions survived the original suite but failed
checks anchored to the supplied data. This complements PatchDiff's use of
behavioural comparisons to expose errors in patches that pass existing
tests~\cite{wang2026solved}. The small change in mutation scores also matters.
Although mutant detection correlates with real-fault detection in prior
work~\cite{just2014mutants}, these scores barely distinguished our suites'
ability to expose agent shortcuts. An input-specific substitution can remain
inactive on the inputs used by a test suite. Evaluating planted shortcuts
alongside mutants, and checking the requested analysis on the supplied data,
can therefore provide complementary evidence of test adequacy.

The disclosure results concern what reaches the user. A reader relying on
the final answer could miss a shortcut mentioned during implementation.
Review brought those qualifications back into view, but an admission does
not repair the analysis. An omission also does not establish deliberate
concealment. Final reports should identify departures from the requested
method and the results they affect. Review can help surface these issues;
assessing correctness still requires independent checks of the code.

In practice, evaluators can combine checks on new inputs with checks that
transform the supplied data, inspect constants or branches that recognise
the dataset, and request a review before accepting the final summary. These
checks guide manual inspection: a constant may be legitimate, and a review
admission still leaves the analysis to be corrected. If tests conflict with the requested
analysis, the discrepancy should be reported rather than resolved by
substituting the expected answer.

The evidence has several limits. Most identified gaming arose under
deliberately conflicting tests. Uneven task and model coverage and incomplete
labels leave ordinary gaming prevalence and the effect of hiding tests
unresolved. We manually assigned labels and disclosure codes. They
remain under author review, without recorded independent double-coding.
Anchored checks were developed after observing pilot shortcuts; their
success therefore shows detection of these collected cases, not performance
on a prospectively held-out set. Retrospective evaluator changes further
limit claims about unseen shortcuts, and
repeated tasks limit generalisation from run-level statistics. Neutral review
preceded direct review in the same session, preventing a clean comparison
of question wording.

The next evaluation should freeze the checks before collecting fresh runs
and have two reviewers independently classify every implementation, including
those that pass. Balancing tasks, models, and prompts across visibility
conditions would make comparisons more informative. Neutral and direct
review questions could be assigned to separate sessions branched from the
same saved implementation context, separating wording from order effects.
These changes would test whether the detection and disclosure patterns
extend beyond the collected cases.
